(1) 经典运动满足 \(p_{\rm cl}(x)=\sqrt{2m[E-a|x|^p]}\)，转向点为 \(x_\pm=\pm x_t\)，\(x_t=(E/a)^{1/p}\)。闭合相空间作用量
\[
J(E)=\oint p_{\rm cl}\,dx
=4\int_0^{x_t}\sqrt{2m(E-ax^p)}\,dx.
\]
作用量量纲为动量乘长度，与 ℏ 相同。

(2) 令 \(x=x_ty\) 后
\[
J(E)=4\sqrt{2m}\,E^{1/2+1/p}a^{-1/p}
\int_0^1\sqrt{1-y^p}\,dy.
\]
再用 \(t=y^p\)，积分为
\[
\frac1p\int_0^1t^{1/p-1}(1-t)^{1/2}dt
=\frac1p B(1/p,3/2).
\]
因此
\[
J(E)=\frac{4\sqrt{2m}}p a^{-1/p}B(1/p,3/2)
E^{(p+2)/(2p)}.
\]

(3) 两端平滑转向点的 WKB 条件是 \(J(E_n)=2\pi\hbar(n+1/2)\)，\(n=0,1,\dots\)。解得
\[
E_n^{\rm WKB}=
\left[\frac{\pi\hbar p\,a^{1/p}(n+1/2)}
{2\sqrt{2m}B(1/p,3/2)}\right]^{2p/(p+2)}.
\]
这明确给出 \(E_n\propto(n+1/2)^{2p/(p+2)}\)，不是仅凭量纲猜测的常数因子；高 n 时 WKB 近似改善。

(4) 对 p=2，令 \(a=m\omega^2/2\)，Beta 函数 \(B(1/2,3/2)=\pi/2\)，代入得到 \(E_n=\hbar\omega(n+1/2)\)，恰与谐振子精确谱吻合。对 p→∞，势的转向点迅速接近硬壁，若简单地把光滑转向点的 Maslov 指数仍固定为 1/2，可能无法恢复无限深方阱的低阶 \(n^2\) 精确量子数偏移；应先选择硬壁边界条件再量子化。这不影响高 n 的平方能级标度。
